\documentclass{beamer}

\newcommand{\week}{Federal Reserve 4 of 5}

\title{Fed Indpendence, FOMC Meetings, and More History}
\subtitle{Mishkin Chapter 16}
\author{Econ 357}
\date{\week}

% Reference the shared preamble
\input{preamble}

% Set the footer -- change 
\setbeamertemplate{footline}{
  \leavevmode%
  \vspace{2ex}
  \hbox{%
    % Left box: Econ 457
    \begin{beamercolorbox}[wd=.4\paperwidth,ht=2.5ex,dp=1ex,left]{author in head/foot}%
      \hspace{1em}Econ 357
    \end{beamercolorbox}%
    % Middle box: Week
    \begin{beamercolorbox}[wd=.2\paperwidth,ht=2.5ex,dp=1ex,center]{date in head/foot}%
      \centering\week
    \end{beamercolorbox}%
    % Right box: Slide numbers
    \begin{beamercolorbox}[wd=.4\paperwidth,ht=2.5ex,dp=1ex,center]{date in head/foot}%
      \hfill\insertframenumber{} 
    \end{beamercolorbox}%
  }%
  \vskip0pt%
}

\begin{document}

\frame{\titlepage}

\begin{frame}
    \frametitle{Outline for Lectures on the Federal Reserve}

    \begin{enumerate}
        \item History / Structure
        \item Money Supply Process
        \item Dual Mandate
        \item \fbox{Fed Independence, FOMC Meetings, and More History}
        \item Extraordinary Tools: Interest Rates and Balance Sheet
    \end{enumerate}
    \vspace{1em}
    Reference: Mishkin Chapters: 13, 14, 15, 16

\end{frame}

\begin{frame}
    \frametitle{Outline for Today's Lecture}
    
    \begin{enumerate}
        \item Fed Independence
        \item FOMC Meetings
        \item More History
        \begin{itemize}
            \item Greenspan and the productivity boom of the 1990s
            \item Yellen and the oil shock of 2011
            \item Powell and the COVID shock of 2020
        \end{itemize}
    \end{enumerate}

\end{frame}

\begin{frame}[t]
    \frametitle{1. Fed Independence}
    \framesubtitle{}
    
    Is the Fed Independent?   It’s complicated!
    \vspace{1em}
    
    Considerations:
    \begin{itemize}
        \item Fed Chairs are appointed by the President.  Sometimes across party lines, but not always.
        \item Fed structure is set up to ensure some stability around elections (14 year terms, Regional Banks, etc)
        \item Fed goals are set by Congress: “maximum employment and price stability”
        \item Congress could always change the Fed’s mandate or take away the Fed’s authority
    \end{itemize}
    
\end{frame}

\begin{frame}[t]
    \frametitle{1. Fed Independence}
    \framesubtitle{}
    
    Considerations (continued):
    \begin{itemize}
        \item Some economists argue that the Fed should be independent to solve principal-agent problems.
        \item Specifically, politicians have short term incentives to boost growth (which leads to inflation).   Delegating authority to the Fed can lead to a better outcome with regards to inflation, as the Fed may be more willing to slow the economy in the near term.
        \item On the other hand, accountability is important.  What should we do when the Fed makes mistakes?   How should the Fed be held accountable?
    \end{itemize}

\end{frame}

\begin{frame}[t]
    \frametitle{1. Fed Independence}
    \framesubtitle{}

    History:
    \begin{itemize}
        \item Nixon put pressure on Burns ahead of the 1974 election. (captured on tape!)   Burns failed to slow economic growth in the 1970s, which led to inflation.
        \item Volcker raised interest rates and caused a recession.  This was politically unpopular, but Reagan reappointed him anyways.
        \item Greenspan had a reputation as a political operative.   Initially appointed by Reagan, reappointed by Bush, Clinton, and G.W. Bush.
        \item Bernanke was appointed by Bush, reappointed by Obama.
        \item Yellen was viewed as a Democrat.  Was appointed by Obama.   Was not reappointed by Trump.
        \item Powell was appointed by Trump.  Then was criticized by Trump in 2019 and investigated by Trump's DOJ in 2025.
    \end{itemize}
    
\end{frame}


\begin{frame}[t]
    \frametitle{1. Fed Independence}
    \framesubtitle{}

    Conclusion….  Is the Fed independent?
    \vspace{1em}

    Make up your own mind.   
    \vspace{1em}
    
    Can be seen from both sides.   
    \vspace{1em}
    
    For better or for worse.
    \vspace{1em}


\end{frame}

\begin{frame}[t]
    \frametitle{2. FOMC Meetings}
    \framesubtitle{}

\begin{itemize}
    \item The FOMC meets every six weeks (eight times per year).
    \item Two day meeting.   
    \item Announcement made on the second day (usually a Wednesday).   \item Announcements include:
    \begin{itemize}
        \item Statement (contains decision on interest rates)
        \item Forecasts (released once per quarter, includes GDP, unemployment, inflation and rates forecast.  Known as the “dot plot”)
        \item Press conference by the Fed Chair
    \end{itemize}
\end{itemize}

\end{frame}

\begin{frame}[t]
    \frametitle{2. FOMC Meetings}
    \framesubtitle{Forecasts}

    \centering
    \includegraphics[width=\textwidth]{Econ357/figures/04_fed_4_forecast.png}

\end{frame}

\begin{frame}[t]
    \frametitle{2. FOMC Meetings}
    \framesubtitle{Dot Plot}

    \centering
    \includegraphics[width=0.8\textwidth]{Econ357/figures/04_fed_4_dots.png}

\end{frame}

\begin{frame}[t]
    \frametitle{2. FOMC Meetings}
    \framesubtitle{}

    The Chair gives a press conference after the announcement
    \vspace{1em}

    \url{https://www.federalreserve.gov}

\end{frame}

\begin{frame}[t]
    \frametitle{2. FOMC Meetings}
    \framesubtitle{}

    The FOMC announcement often causes markets to adjust:
    \vspace{1em}

    \url{https://www.google.com/url?sa=t&source=web&rct=j&opi=89978449&url=https://www.cnbc.com/video/2024/03/20/short-the-days-of-powell-shocking-the-market-are-over-says-western-assetas-john-bellows.html&ved=2ahUKEwiWsaC-g5uTAxXjk2oFHcWAENEQtwJ6BAgQEAI&usg=AOvVaw2VUv7UpcRHIhYm6lF9NPyd}

\end{frame}

\begin{frame}[t]
    \frametitle{3. More History}
    \framesubtitle{}

    What does the Fed do when there is tension between the two parts of the mandate?   Say, for example, that inflation is too high AND unemployment is also too high?  
    \vspace{1em}
    
    The Taylor Rule provides some sense of the trade-offs.
    \vspace{1em}
    
    It's also useful to look at historical precedents.

\end{frame}

\begin{frame}[t]
    \frametitle{3. More History}
    \framesubtitle{Three Episodes from History}

    \begin{enumerate}
        \item In the 1990s, the Greenspan Fed did not raise interest rates, even though inflation was a bit too high.   Greenspan  believed that inflation would come down due to technology and an increase in productivity.   In retrospect, the Fed made the \textbf{right} decision.
        \item In 2011, the Fed did not raise interest rates even though inflation was rising.   The Fed, and especially then-Vice Chair Janet Yellen, thought that a still-high unemployment rate would bring down inflation on its own.   In retrospect, the Fed made the \textbf{right} decision.
        \item In 2022, the Powell Fed was slow to raise interest rates even though inflation was rising.   The Fed thought the inflation would be transitory.   In retrospect, the Fed probably made the \textbf{wrong} decision.
    \end{enumerate}

\end{frame}

\begin{frame}[t]
    \frametitle{3. More History}
    \framesubtitle{1990s}

        \begin{columns}
            \begin{column}{0.5\textwidth}
            \textbf{Alan Greenspan}
        \begin{itemize}
            \item Proponent of laissez-faire economics (Ayn Rand)
            \item Oversaw the disinflationary boom of the 1990s.   Known as “The Maestro”
            \item Was a proponent of financial deregulation in the 1990s and 2000s, which may have contributed to the Financial Crisis of 2008
        \end{itemize}
        \end{column}
        
        \begin{column}{0.5\textwidth}
            \centering
            \includegraphics[width=0.5\textwidth]{Econ357/figures/04_fed_4_greenspan.png}
        \end{column}
    
    \end{columns}

\end{frame}

\begin{frame}[t]
    \frametitle{3. More History}
    \framesubtitle{1990s}

    \centering
    \includegraphics[width=\textwidth]{Econ357/figures/04_fed_4_1990sue.png}

\end{frame}

\begin{frame}[t]
    \frametitle{3. More History}
    \framesubtitle{1990s}

    \centering
    \includegraphics[width=\textwidth]{Econ357/figures/04_fed_4_1990scpi.png}

\end{frame}

\begin{frame}[t]
    \frametitle{3. More History}
    \framesubtitle{1990s}

    \centering
    \includegraphics[width=\textwidth]{Econ357/figures/04_fed_4_1990sff.png}

\end{frame}

\begin{frame}[t]
    \frametitle{3. More History}
    \framesubtitle{1990s}

    In retrospect, Greenspan made the right decision to not raise rates in the early 1990s.   Appropriately, Greenspan is often credited for correctly forecasting the disinflation in the 1990s, and for prioritizing the labor market during this period.
    \vspace{1em}
    
    Criticisms of Greenspan:
    \begin{itemize}
        \item In 1996 Greenspan called out the “irrational exuberance” of the stock market.   But, he did nothing to slow down the bubble that formed over the next 4 years.
        \item Greenspan supported deregulation, and did too little to reign in bank lending in the 2000s.
    \end{itemize}


\end{frame}

\begin{frame}[t]
    \frametitle{3. More History}
    \framesubtitle{2010s}

        \begin{columns}
            \begin{column}{0.5\textwidth}
            \textbf{Ben Bernanke}
            \begin{itemize}
                \item Professor of Economic History.   Won the Nobel Prize in 2022 for his work on the Great Depression.
                \item Guided the Fed through the 2008 Financial Crisis.
            \end{itemize}
        \end{column}
        
        \begin{column}{0.5\textwidth}
            \centering
            \includegraphics[width=0.5\textwidth]{Econ357/figures/04_fed_4_bernanke.png}
        \end{column}
    
    \end{columns}

\end{frame}

\begin{frame}[t]
    \frametitle{3. More History}
    \framesubtitle{2010s}

        \begin{columns}
            \begin{column}{0.5\textwidth}
            \textbf{Janet Yellen}
         \begin{itemize}
             \item Professor of Economics, focus on labor economics
             \item Long history of public service (SF Fed, Vice Chair at the Fed)
             \item Treasury Secretary under President Biden, 2021-2025
         \end{itemize}
        

        \end{column}
        
        \begin{column}{0.5\textwidth}
            \centering
            \includegraphics[width=0.5\textwidth]{Econ357/figures/04_fed_4_yellen.png}
        \end{column}
    
    \end{columns}

\end{frame}

\begin{frame}[t]
    \frametitle{3. More History}
    \framesubtitle{2010s}

    \centering
    \includegraphics[width=\textwidth]{Econ357/figures/04_fed_4_2011cpi.png}

\end{frame}

\begin{frame}[t]
    \frametitle{3. More History}
    \framesubtitle{2010s}

    \centering
    \includegraphics[width=\textwidth]{Econ357/figures/04_fed_4_2011ue.png}

\end{frame}

\begin{frame}[t]
    \frametitle{3. More History}
    \framesubtitle{2010s}

    In retrospect, the Fed made the correct decision to not raise rates in 2011.   The increase in inflation was not sustained, due in large part to the slack in the labor market.
    \vspace{1em}

    Not only did inflation fall in the following years, it actually fell to below target, so that the Fed was forced to try a number of things to further ease policy in subsequent years.

\end{frame}

\begin{frame}[t]
    \frametitle{3. More History}
    \framesubtitle{2021-2023}

        \begin{columns}
            \begin{column}{0.5\textwidth}
            \textbf{Jerome Powell}
        \begin{itemize}
            \item Lawyer (not an economist)
            \item Appointed by President Trump as a “low interest rate Republican”.
            \item Criticized by Trump for raising interest rates in 2018 and for not cutting rates in 2025.
            \item Oversaw an aggressive response to COVID
            \item Struggled to respond to post-COVID inflation.
        \end{itemize}
        \end{column}
        
        \begin{column}{0.5\textwidth}
            \centering
            \includegraphics[width=0.5\textwidth]{Econ357/figures/04_fed_4_powell.png}
        \end{column}
    
    \end{columns}

\end{frame}

\begin{frame}[t]
    \frametitle{3. More History}
    \framesubtitle{2021-2023}

    \centering
    \includegraphics[width=\textwidth]{Econ357/figures/04_fed_4_2022ue.png}

\end{frame}

\begin{frame}[t]
    \frametitle{3. More History}
    \framesubtitle{2021-2023}

    \centering
    \includegraphics[width=\textwidth]{Econ357/figures/04_fed_4_2022ff.png}

\end{frame}

\begin{frame}[t]
    \frametitle{3. More History}
    \framesubtitle{2021-2023}

    In retrospect, the Fed made a mistake by not raising rates in 2021.   The Fed believed inflation would be transitory, yet this was wrong and inflation remained high even into 2022.
    \vspace{1em}
    
    By the time the Fed did raise rates in 2022, they were forced to raise rates very rapidly (3475 bps of hikes in total, including 4x 75 bps moves).

\end{frame}

\end{document}